\documentclass{beamer}
\usetheme{Madrid}

\usepackage[utf8]{inputenc} % solo si usas pdflatex
\usepackage[T1]{fontenc}    % solo si usas pdflatex
\usepackage[spanish,es-noshorthands]{babel}

\usepackage{amsmath,amssymb}
\usepackage{microtype}
\usepackage{cancel}         % si lo usas
\usepackage{fontawesome}    % si lo usas
% \usepackage{xypic}        % solo si lo usas
% \usepackage{multicol}     % normalmente no; usar columns de beamer
% \usepackage{amsthm}       % solo si defines teoremas propios
% \usepackage{scrextend}    % solo si lo usas

\usepackage{graphicx}
\graphicspath{{images/}}

\usepackage{booktabs}
\usepackage{tabularx}
\usepackage{array}
\usepackage{ragged2e}

\hypersetup{pdfborder={0 0 0}}

\title{Introducción a Agile}
\author{Rodrigo Blanco Betes. Alfa Formación}
\date{Enero del 2026}

\begin{document}

\begin{frame}
\titlepage
\end{frame}

\begin{frame}{Contenido}
\tableofcontents
\end{frame}

\section{Agile y Scrum}

\begin{frame}{Del enfoque tradicional a Agile}
\begin{columns}[T,onlytextwidth]
\column{0.49\textwidth}
\textbf{Problema del enfoque tradicional}
\begin{itemize}
    \item Planificación cerrada desde el inicio.
    \item Cambios costosos y lentos de incorporar.
    \item Entrega de valor al final del proyecto.
    \item Riesgo de descubrir tarde errores o necesidades reales.
\end{itemize}

\column{0.49\textwidth}
\textbf{Principios de Agile}
\begin{itemize}
    \item Trabajo iterativo e incremental.
    \item Adaptación continua al cambio.
    \item Colaboración con cliente y equipo.
    \item Entregas frecuentes con valor real.
\end{itemize}
\end{columns}
\end{frame}

\begin{frame}{Modelo tradicional vs Agile}
\vspace{0.4cm}
\small
\begin{tabularx}{\textwidth}{>{\RaggedRight\arraybackslash}p{0.25\textwidth} >{\RaggedRight\arraybackslash}X >{\RaggedRight\arraybackslash}X}
\toprule
\textbf{Aspecto} & \textbf{Tradicional} & \textbf{Agile} \\
\midrule
Planificación & Se define al principio & Se revisa en cada iteración \\
Entrega & Al final & Frecuente y progresiva \\
Cambios & Se intentan evitar & Se gestionan como parte natural \\
Éxito & Cumplir el plan & Aportar valor y aprender rápido \\
\bottomrule
\end{tabularx}
\end{frame}

\begin{frame}{¿Qué es Scrum?}
\begin{itemize}
    \item Scrum es un marco de trabajo Agile para desarrollar proyectos complejos de forma iterativa.
    \item El trabajo se organiza en \textbf{sprints}: periodos cortos y fijos en los que se entrega un incremento usable.
    \item Se apoya en tres ideas clave:
    \begin{itemize}
        \item \textbf{Transparencia}: todos saben qué se está haciendo y con qué prioridad.
        \item \textbf{Inspección}: el equipo revisa con frecuencia el avance y el resultado.
        \item \textbf{Adaptación}: se ajusta el rumbo en cuanto aparece nueva información.
    \end{itemize}
    \item Su objetivo no es solo \textit{hacer tareas}, sino \textbf{entregar valor con control del proceso}.
\end{itemize}
\end{frame}

\begin{frame}{Roles en Scrum y relación entre ellos}
\begin{columns}[T,onlytextwidth]
\column{0.32\textwidth}
\textbf{Product Owner}
\begin{itemize}
    \item Define prioridades.
    \item Maximiza el valor.
    \item Gestiona el Product Backlog.
\end{itemize}

\column{0.32\textwidth}
\textbf{Scrum Master}
\begin{itemize}
    \item Facilita el marco Scrum.
    \item Elimina impedimentos.
    \item Ayuda a mejorar la forma de trabajar.
\end{itemize}

\column{0.32\textwidth}
\textbf{Equipo de Desarrollo}
\begin{itemize}
    \item Organiza cómo realizar el trabajo.
    \item Construye el incremento.
    \item Colabora de forma multidisciplinar.
\end{itemize}
\end{columns}

\vspace{0.35cm}
\begin{block}{Relación entre roles}
El \textbf{Product Owner} marca el \textit{qué} y el \textit{por qué}; el \textbf{equipo} decide el \textit{cómo}; el \textbf{Scrum Master} cuida el proceso para que todo fluya.
\end{block}
\end{frame}

\begin{frame}{Ceremonias de Scrum para el control del proceso}
\small
\begin{tabularx}{\textwidth}{>{\RaggedRight\arraybackslash}p{0.22\textwidth} >{\RaggedRight\arraybackslash}p{0.22\textwidth} X}
\toprule
\textbf{Ceremonia} & \textbf{Objetivo} & \textbf{Cómo ayuda al control} \\
\midrule
Sprint Planning & Decidir qué se hará en el sprint & Alinea prioridades, capacidad y compromiso \\
Daily Scrum & Revisar avance diario & Detecta bloqueos y desviaciones con rapidez \\
Sprint Review & Mostrar el incremento & Valida resultado y recoge feedback real \\
Sprint Retrospective & Mejorar la forma de trabajar & Identifica aprendizajes y acciones de mejora \\
\bottomrule
\end{tabularx}

\vspace{0.25cm}
\normalsize
\textbf{Idea clave:} las ceremonias no son reuniones por rutina; son puntos de inspección y adaptación que mantienen el proyecto bajo control.
\end{frame}

\begin{frame}{Ejemplo práctico de uso de Scrum}
\textbf{Proyecto:} lanzar una nueva experiencia digital de reservas.

\vspace{0.2cm}
\begin{enumerate}
    \item El Product Owner prioriza funcionalidades: búsqueda, pago y confirmación.
    \item El equipo selecciona en Sprint Planning lo que puede completar en dos semanas.
    \item Cada día revisa avances y bloqueos en la Daily.
    \item Al final del sprint enseña una versión funcional en la Review.
    \item En la Retrospective ajusta la forma de colaborar para el siguiente sprint.
\end{enumerate}

\vspace{0.25cm}
\begin{alertblock}{Resultado}
Menos riesgo, mayor visibilidad, capacidad de corregir antes y entregas útiles desde fases tempranas.
\end{alertblock}
\end{frame}

\section{Design Thinking}

\begin{frame}{Design Thinking: innovación centrada en las personas}
\begin{itemize}
    \item Es una metodología para encontrar soluciones innovadoras a partir de la empatía, la creatividad y la experimentación.
    \item Parte de una pregunta básica: \textbf{¿qué necesita realmente la persona usuaria?}
    \item No busca solo ideas originales, sino ideas \textbf{deseables, viables y factibles}.
\end{itemize}

\vspace{0.3cm}
\textbf{Etapas principales}
\begin{enumerate}
    \item Empatizar
    \item Definir
    \item Idear
    \item Prototipar
    \item Testear
\end{enumerate}
\end{frame}

\begin{frame}{Etapas de Design Thinking}
\small
\begin{tabularx}{\textwidth}{>{\RaggedRight\arraybackslash}p{0.18\textwidth} X >{\RaggedRight\arraybackslash}p{0.30\textwidth}}
\toprule
\textbf{Etapa} & \textbf{Propósito} & \textbf{Pregunta guía} \\
\midrule
Empatizar & Comprender personas, contexto y necesidades & ¿Qué vive, siente y necesita el usuario? \\
Definir & Concretar el problema correcto & ¿Qué reto merece ser resuelto? \\
Idear & Generar alternativas sin cerrar pronto la solución & ¿Qué opciones podríamos explorar? \\
Prototipar & Dar forma rápida a las ideas & ¿Cómo convertir la idea en algo visible? \\
Testear & Aprender con feedback real & ¿Qué funciona y qué debemos mejorar? \\
\bottomrule
\end{tabularx}
\end{frame}

\begin{frame}{Técnicas útiles en cada fase}
\begin{columns}[T,onlytextwidth]
\column{0.48\textwidth}
\textbf{Comprender el problema}
\begin{itemize}
    \item Entrevistas
    \item Observación
    \item Mapa de empatía
    \item \textit{Customer journey}
\end{itemize}

\textbf{Definir e idear}
\begin{itemize}
    \item \textit{How might we...}
    \item Brainstorming
    \item Brainwriting
    \item Priorización impacto/esfuerzo
\end{itemize}

\column{0.48\textwidth}
\textbf{Prototipar}
\begin{itemize}
    \item Bocetos
    \item Storyboards
    \item Maquetas rápidas
    \item Role play
\end{itemize}

\textbf{Testear}
\begin{itemize}
    \item Pruebas con usuarios
    \item Entrevistas de validación
    \item Observación de uso
    \item Iteración sobre hallazgos
\end{itemize}
\end{columns}
\end{frame}

\begin{frame}{Cómo se complementan Design Thinking y Scrum}
\begin{itemize}
    \item \textbf{Design Thinking} ayuda a descubrir el problema correcto y a diseñar soluciones con sentido.
    \item \textbf{Scrum} ayuda a convertir esas soluciones en entregas reales, iterativas y controladas.
\end{itemize}

\vspace{0.3cm}
\begin{block}{Secuencia habitual}
Primero exploramos necesidades, generamos ideas y validamos conceptos. Después organizamos la construcción mediante backlog, sprints, revisiones y mejora continua.
\end{block}

\vspace{0.25cm}
\textbf{Resultado:} más creatividad al inicio y más disciplina en la ejecución.
\end{frame}

\section{Storytelling}

\begin{frame}{Storytelling: el arte de transmitir un mensaje con éxito}
\begin{itemize}
    \item Contar historias permite convertir datos o ideas en mensajes memorables.
    \item Una buena historia genera atención, comprensión y conexión emocional.
    \item No consiste en adornar el mensaje, sino en \textbf{darle estructura, sentido y fuerza}.
\end{itemize}

\vspace{0.3cm}
\textbf{¿Para qué sirve?}
\begin{itemize}
    \item Convencer.
    \item Inspirar.
    \item Explicar cambios o proyectos.
    \item Facilitar que el público recuerde el mensaje principal.
\end{itemize}
\end{frame}

\begin{frame}{Estructura y elementos de una buena historia}
\begin{columns}[T,onlytextwidth]
\column{0.48\textwidth}
\textbf{Estructura básica}
\begin{enumerate}
    \item Contexto inicial
    \item Conflicto o reto
    \item Búsqueda de solución
    \item Resultado o aprendizaje
\end{enumerate}

\column{0.48\textwidth}
\textbf{Elementos clave}
\begin{itemize}
    \item Personaje o protagonista
    \item Objetivo claro
    \item Obstáculo o tensión
    \item Cambio final
    \item Mensaje central
\end{itemize}
\end{columns}

\vspace{0.25cm}
\begin{block}{Regla práctica}
Si la audiencia puede resumir la idea en una frase al terminar, la historia ha funcionado.
\end{block}
\end{frame}

\begin{frame}{Consejos y ejemplo breve}
\textbf{Consejos}
\begin{itemize}
    \item Empieza con una situación reconocible.
    \item Evita saturar con datos sin hilo conductor.
    \item Da protagonismo al cambio: antes, durante y después.
    \item Cierra con una idea clara o llamada a la acción.
\end{itemize}

\vspace{0.25cm}
\textbf{Ejemplo}
\begin{quote}
``El equipo recibía muchas incidencias porque el proceso era lento y confuso. Tras escuchar a los usuarios, rediseñamos los pasos críticos y probamos una versión simple. El tiempo de gestión bajó y la satisfacción subió. La lección fue clara: entender bien el problema cambia la calidad de la solución.''
\end{quote}
\end{frame}

\begin{frame}{Puesta en práctica}
\textbf{Dinámica sugerida}
\begin{enumerate}
    \item Elegir un caso real o un proyecto cercano.
    \item Describir el problema inicial en una frase.
    \item Construir una mini historia con esta plantilla:
    \begin{itemize}
        \item \textbf{Situación}: qué estaba ocurriendo.
        \item \textbf{Reto}: qué dificultad apareció.
        \item \textbf{Acción}: qué hicimos.
        \item \textbf{Resultado}: qué cambió.
    \end{itemize}
    \item Contarla en 60-90 segundos.
    \item Recibir feedback sobre claridad, impacto y memorabilidad.
\end{enumerate}
\end{frame}

\begin{frame}{Ideas clave para cerrar}
\begin{itemize}
    \item \textbf{Agile y Scrum} ayudan a gestionar proyectos con flexibilidad, foco y control.
    \item \textbf{Design Thinking} ayuda a descubrir problemas relevantes y generar soluciones innovadoras.
    \item \textbf{Storytelling} ayuda a comunicar ideas de manera clara, persuasiva y memorable.
\end{itemize}

\vspace{0.35cm}
\centering
\Large
Pensar bien, construir bien y comunicar bien.
\end{frame}

\end{document}